\documentclass[11pt]{article}
% Shared preamble of the RMC kernel write-ups (\input{preamble} after \documentclass).
\usepackage[margin=1in]{geometry}
\usepackage{amsmath,amssymb,bm}
\usepackage{graphicx}
\usepackage{booktabs}
\usepackage{hyperref}

\newcommand{\dd}{\,\mathrm{d}}
\newcommand{\Ein}{E}
\newcommand{\Eout}{E'}
\newcommand{\Ecm}{E_{c}}
\newcommand{\mucm}{\mu_{c}}
\newcommand{\mulab}{\mu_0}

\newcommand{\acespec}[1]{J.~L.~Conlin and P.~Romano,
\href{https://mcnp-green.lanl.gov/pdf_files/TechReport_2019_LANL_LA-UR-19-29016_ConlinRomano.pdf}
{\emph{A Compact ENDF (ACE) Format Specification}}, LA-UR-19-29016 (2019),
Sec.~#1}
\newcommand{\ENDFspec}[1]{\href{https://www.nndc.bnl.gov/endfdocs/ENDF-102-2023.pdf}
{\emph{ENDF-6 Formats Manual}} (2023), #1}

% Common closing section: how the verification figures are produced
\newcommand{\verificationmethod}{%
Each figure is produced by \texttt{verification/verify\_laws.py}. For an incident
energy $\Ein$ along $+z$, $10^6$ collisions are sampled with MC/DC's own reaction
sampler (\texttt{sample\_elastic\_scattering}, \texttt{sample\_inelastic\_scattering}
or \texttt{sample\_fission}), and every emitted neutron is histogrammed in
$48 \times 24$ bins of $(\Eout, \mulab)$, with $\mulab$ the lab cosine to $+z$. The
inverted kernel used by RMC is integrated over the same bins (Gauss--Legendre panels;
two-body lines are split where $\mulab^*(\Eout)$ crosses a cosine edge), multiplied by
the number of collisions, and compared with a Pearson-type $\chi^2$ diagnostic
over bins with at least 5 expected counts (the rest pooled into one bin).
The existing tables list the script's $\chi^2/\mathrm{dof}$ and nominal $p$-value,
the emitted neutrons per collision (sampled vs.\ the integral of the
inverted kernel), and the largest relative change of a bin integral when the
quadrature panels are doubled. Left panels: outgoing-energy marginal; middle:
lab-cosine marginal; right: per-bin $z = (\text{sampled} - \text{inverted}) /
\sqrt{\text{inverted}}$.

\noindent These figures check agreement with the MC/DC sampler, not independent
physical validity. Their existing $p$-values are not calibrated goodness-of-fit
probabilities: the script uses $B$ degrees of freedom, whereas a fixed-total
multinomial histogram with $B$ categories and no fitted parameters uses $B-1$.
Random multiplicities and spectrum assignments can also change the collision-level
covariance from independent Poisson counts; $z$ is then only a count diagnostic.
See \href{https://docs.scipy.org/doc/scipy/reference/generated/scipy.stats.chisquare.html}
{the Pearson test's sampling assumptions}.
Energy bins use the sample extrema, so unsampled tails are excluded and the
reported inverted yield is an integral over that observed window. A formal test
needs independently specified bins, explicit tail categories, and a covariance or
resampling model appropriate to independent collisions. Zero expected density
with observed emissions must be reported rather than discarded. Doubling quadrature
panels gives a convergence indicator, not a rigorous error bound. The stored
figures and tables have not been regenerated for these qualifications.}


\newcommand{\tpsi}{\tilde\psi}
\newcommand{\teps}{\tilde\epsilon}
\newcommand{\Sbar}{\bar S}
\newcommand{\E}{\mathbb{E}}

\title{Allocating correction-pass histories between residual terms}
\date{}

\begin{document}
\maketitle

\noindent A correction pass (\texttt{method/rmc\_math.tex}) samples the full residual as
separate terms,
\begin{equation}
  r = \underbrace{(r_c + r_e)}_{i=1} + \underbrace{r_s}_{i=2} + \underbrace{r_f}_{i=3},
  \qquad r_s = T_s - \Sbar_s,\quad r_f = T_f - \Sbar_f ,
\end{equation}
each with its own proposal and its own number of histories $N_i$,
$\sum_i N_i = N$. Until now the split was fixed: $N_1 = 0.9N$
(\texttt{collision\_fraction}), and the rest went to $r_s$ and $r_f$. This note derives
the split that minimizes the variance of a pass and describes its implementation
(\texttt{correction\_allocation="mass"}, \texttt{driver.allocate\_histories}).

\section{The optimal split}

A pass estimates a linear response $R$ of $\teps = L^{-1}r$, for example a bin average,
as a sum of independent term estimates,
\begin{equation}
  \hat R = \sum_i \hat R_i,\qquad
  \hat R_i = \frac{1}{N_i}\sum_{k=1}^{N_i} Y_{ik},\qquad
  \operatorname{Var}\hat R = \sum_i \frac{\sigma_i^2}{N_i},
\end{equation}
with $Y_{ik}$ the response of history $k$ of term $i$ (its source weight times the
transported response) and $\sigma_i^2$ its per-history variance. Each $\hat R_i$ is
unbiased for any $N_i\ge1$, so the split affects only the variance. Minimizing
$\sum_i\sigma_i^2/N_i$ subject to $\sum_iN_i=N$ (a Lagrange multiplier) gives
\begin{equation}
  N_i = N\,\frac{\sigma_i}{\sum_k\sigma_k},\qquad
  \operatorname{Var}\hat R = \frac{\big(\sum_i\sigma_i\big)^2}{N},
  \label{eq:neyman}
\end{equation}
the optimal allocation of stratified sampling [1, 2, 3].

\paragraph{Why the $L^1$ masses.} Each term is sampled with a proposal proportional to
its own $|r_i|$, by bin (the collision and edge masses, \texttt{method/rmc\_math.tex}; the
correction masses of the integrated sampler). With such a proposal every source weight
has magnitude $\|r_i\|_1$, exactly when $|r_i|$ is constant in each bin and nearly so
otherwise, so
\begin{equation}
  \sigma_i \approx \|r_i\|_1\,\rho_i,
\end{equation}
with $\rho_i^2$ the second moment of the response to one unit-weight history of term
$i$. The terms start in the same phase space and share the transport, so their
$\rho_i$ are comparable, and \eqref{eq:neyman} becomes
\begin{equation}
  N_i \propto \|r_i\|_1 ,
\end{equation}
which is what \texttt{correction\_allocation="mass"} uses.

\paragraph{What the fixed split costs.} With fractions $f_i$ and masses $m_i =
\|r_i\|_1$, the source second moment is $\sum_i m_i^2/f_i$, against
$(\sum_i m_i)^2$ for the mass split, a ratio of
\begin{equation}
  \frac{\sum_i m_i^2/f_i}{(\sum_i m_i)^2}\ \ge 1 ,
\end{equation}
with equality only when $f_i = m_i/\sum_k m_k$. Near the collision-only fixed point the
binned residual $r_c + r_e$ is small, by construction of the fixed point, and $r_s$
holds most of the mass. With $m_s$ at 95\% of the total and $f_s = 0.1$ the ratio is
$0.05^2/0.9 + 0.95^2/0.1 \approx 9$. A 0D prototype of these samplers measured the same
effect, about $10\times$ excess variance per pass at the fixed point [4].

\section{Implementation}

\begin{itemize}
  \item $\|r_c + r_e\|_1$ is exact: the sum of the collision and edge sampling masses
  (\texttt{Residual.collision\_edge\_norm}).
  \item $\|r_s\|_1$ and $\|r_f\|_1$ come from the integrated sampler's deterministic
  pilot, which already evaluates $|T - \Sbar|$ at $2\times2$ Gauss points per bin
  (\texttt{source.correction\_masses\_l1}). The proposal adds a floor of
  $0.1\,|\Sbar|$ per bin, so that every bin receiving in-scatter can be sampled. That
  floor is not residual mass and is excluded from the allocation; near the fixed point
  it can exceed $\|T-\Sbar\|_1$.
  \item The pilot is an estimate, so its error affects only the efficiency of the
  split, not the bias.
  \item Rounding uses largest remainders, so that $\sum_iN_i=N$ exactly, and gives
  every term with positive mass at least one history.
  \item The histories of each pass are stored in \texttt{RMCResult.correction\_allocation}.
\end{itemize}

\paragraph{The pointwise sampler keeps the fixed split.} Its source weights are far
from $\|r_s\|_1$ in magnitude: on the A\,=\,1 problem its second moment is several
hundred times $\|r_s\|_1^2$ (\texttt{method/scattering\_correction\_samplers.tex}). So
$\sigma_s\not\approx\|r_s\|_1\rho_s$, and the mass rule would give $r_s$ too few
histories. It also lacks the pilot that gives $\|r_s\|_1$. With
\texttt{correction\_allocation="auto"} (the default), the integrated sampler uses
\texttt{"mass"} and the pointwise sampler \texttt{"fixed"}. The general rule
\eqref{eq:neyman} with $\sigma_i$ estimated from a pilot of actual weights would
cover both; it is not implemented.

\section{Measurement}

\texttt{examples/rmc/residual\_allocation/measure.py} uses the A\,=\,1 slowing-down
problem of \texttt{test\_analytic\_slowing\_down\_A1} ($c=1/2$, $G=20$ collided bins,
analytic flux). It runs the collision-only iterations to the fixed point, then 30
correction passes with each split from that same fixed point. Its output is in
\texttt{results.txt}:

\begin{center}
\begin{tabular}{lrrr}
\toprule
 & \texttt{"fixed"} (0.9) & \texttt{"mass"} & ratio \\
\midrule
histories $r_c+r_e$ / $r_s$ / $r_f$ & 18\,900 / 2\,100 / 0 & 1 / 20\,999 / 0 & \\
summed per-bin variance of a pass & $4.37\times10^{-7}$ & $4.05\times10^{-8}$ & 10.8 \\
per-bin variance ratio (median, min--max) & & & 9.8 (4.1--14.6) \\
rms error, fixed point + mean correction & $1.04\times10^{-3}$ & $6.5\times10^{-4}$ & \\
\bottomrule
\end{tabular}
\end{center}
\noindent The two splits agree in mean (largest difference 1.7 standard errors over the
20 bins), so the gain is variance alone. Its size, $10.8\times$, matches the estimate
of Sec.~1 with $m_s$ near 100\% of the mass. Here the cross sections are constant in
energy, so at the collision-only fixed point $r_c$ vanishes, $r_e = 0$ in an infinite
medium, and all of the remaining residual is the in-bin shape of the in-scatter: the
mass split gives $r_c+r_e$ a single history. With energy-dependent cross sections or
faces, $r_c+r_e$ keeps a share in proportion to its mass.

\begin{thebibliography}{9}
\bibitem{neyman} J.~Neyman, ``On the two different aspects of the representative
  method: the method of stratified sampling and the method of purposive selection,''
  \emph{J.\ R.\ Stat.\ Soc.} \textbf{97}(4), 558--625 (1934).
\bibitem{cochran} W.~G.~Cochran, \emph{Sampling Techniques}, 3rd ed., Wiley (1977)
  (optimum allocation in stratified sampling).
\bibitem{owen} A.~B.~Owen, \emph{Monte Carlo Theory, Methods and Examples} (2013),
  chapters on stratification.
\bibitem{prototype} 0D residual-sampling prototype, branch
  \texttt{claude/rmc-residual-sampling-uyqnqi},
  \texttt{examples/rmc/residual\_sampling\_0d/split\_vs\_unsplit*.py} (per-pass
  variance of the fixed and mass-proportional splits at a frozen trial function).
\end{thebibliography}

\end{document}
